\section{ISAC-Assisted Underserved-Region Model}
\label{sec:mathematical-model}

We consider a set of ground UEs $\mathcal{U}$, sensing nodes
$\mathcal{S}$, terrestrial gNBs $\mathcal{M}$, and UAV gNBs
$\mathcal{Q}$. A sensing node may be a terrestrial gNB or a UAV equipped
with ISAC capability. The proposed control loop separates two functions:
the communication measurements determine \emph{which} UEs are
underserved, whereas ISAC observations estimate \emph{where} those UEs
are located. The near-RT RIC fuses both sources to construct an
underserved-region map and determine UAV repositioning targets.

\subsection{System-Level ISAC Sensing Model}

Let $\mathbf{p}_s(t)\in\mathbb{R}^3$ and
$\mathbf{p}_u(t)\in\mathbb{R}^3$ denote the positions of sensing node
$s\in\mathcal{S}$ and UE $u\in\mathcal{U}$ at time $t$. Their sensing
range is
\begin{equation}
 r_{s,u}(t)=\left\|\mathbf{p}_s(t)-\mathbf{p}_u(t)\right\|_2.
 \label{eq:isac-range}
\end{equation}
At the system level, the received echo power is represented by a
monostatic sensing-link abstraction,
\begin{equation}
 P_{s,u}^{\mathrm{rx}}(t)=
 \frac{P_sG_s^2\lambda^2\sigma_u}
 {(4\pi)^3r_{s,u}^4(t)L_s},
 \label{eq:sensing-power}
\end{equation}
where $P_s$, $G_s$, $\lambda$, $\sigma_u$, and $L_s$ are the sensing
transmit power, antenna gain, wavelength, target radar cross section,
and aggregate sensing loss, respectively. The corresponding sensing
signal-to-noise ratio is
\begin{equation}
 \gamma_{s,u}(t)=
 \frac{P_{s,u}^{\mathrm{rx}}(t)}{N_0B_s},
 \label{eq:sensing-snr}
\end{equation}
where $N_0$ and $B_s$ denote the noise spectral density and sensing
bandwidth. This abstraction evaluates sensing availability rather than
waveform-level target estimation.

The probability that sensor $s$ detects UE $u$ is modelled as
\begin{equation}
 P_{s,u}^{\mathrm{det}}(t)=
 \frac{1}
 {1+\exp[-a(\gamma_{s,u}^{\mathrm{dB}}(t)-\gamma_0)]},
 \label{eq:detection-probability}
\end{equation}
where $a$ controls the transition steepness and $\gamma_0$ is the
50\% detection-SNR point. A detection indicator is sampled as
$d_{s,u}(t)\sim\operatorname{Bernoulli}
(P_{s,u}^{\mathrm{det}}(t))$.

For a successful detection, the position estimate is
\begin{equation}
 \widehat{\mathbf{p}}_{s,u}(t)
 =\mathbf{p}_u(t)+\mathbf{e}_{s,u}(t),\qquad
 \mathbf{e}_{s,u}(t)\sim
 \mathcal{N}\!\left(\mathbf{0},
 \sigma_{s,u}^2(t)\mathbf{I}\right),
 \label{eq:sensed-position}
\end{equation}
where the localisation uncertainty decreases with sensing SNR:
\begin{equation}
 \sigma_{s,u}(t)=\operatorname{clip}\!\left[
 \sigma_{\mathrm{ref}}
 10^{-\frac{\gamma_{s,u}^{\mathrm{dB}}(t)-\gamma_{\mathrm{ref}}}{20}},
 \sigma_{\min},\sigma_{\max}\right].
 \label{eq:localisation-uncertainty}
\end{equation}
Here, $\operatorname{clip}[x,l,h]$ limits $x$ to $[l,h]$.
When several sensing nodes detect the same UE, the near-RT RIC forms
an inverse-variance weighted estimate
\begin{equation}
 \widehat{\mathbf{p}}_u(t)=
 \frac{\displaystyle\sum_{s\in\mathcal{S}_u(t)}
 \sigma_{s,u}^{-2}(t)\widehat{\mathbf{p}}_{s,u}(t)}
 {\displaystyle\sum_{s\in\mathcal{S}_u(t)}
 \sigma_{s,u}^{-2}(t)},
 \label{eq:sensing-fusion}
\end{equation}
where $\mathcal{S}_u(t)=\{s:d_{s,u}(t)=1\}$.
The true position $\mathbf{p}_u(t)$ is used only as simulation ground
truth; the controller receives $\widehat{\mathbf{p}}_u(t)$.

\subsection{Communication-Aware Underserved-UE Identification}

At each E2 reporting epoch, let $R_u(t)$, $P_u(t)$, $C_u(t)$,
$D_u(t)$, and $L_u(t)$ denote the serving-cell RSRP, packet-delivery
ratio, throughput, one-way delay, and serving-cell load observed for
UE $u$. The UE is classified as underserved according to
\begin{equation}
 I_u(t)=\mathbf{1}\!\left\{
 R_u(t)<R_{\min}\ \lor\
 P_u(t)<P_{\min}\ \lor\
 C_u(t)<C_{\min}\ \lor\
 D_u(t)>D_{\max}
 \right\}.
 \label{eq:underserved-indicator}
\end{equation}
Only KPIs available in a particular experiment are included in
\eqref{eq:underserved-indicator}; no unmeasured KPI is synthetically
introduced. Consequently, sensing does not itself determine whether
a UE is underserved. It supplies the spatial estimate associated
with the communication-derived indicator $I_u(t)$.

\subsection{Spatial Fusion and Underserved Heatmap}

The horizontal service area is divided into grid cells
$g\in\mathcal{G}$ with centres $\mathbf{c}_g\in\mathbb{R}^2$. Let
$K_h(\cdot)$ be a spatial binning or kernel function with bandwidth
$h$. The sensed UE density and underserved intensity in grid cell
$g$ are
\begin{align}
 \rho_g(t)&=\sum_{u\in\mathcal{U}}
 K_h\!\left(\widehat{\mathbf{p}}_u^{xy}(t)-\mathbf{c}_g\right),
 \label{eq:sensed-density}\\
 H_g(t)&=\sum_{u\in\mathcal{U}} I_u(t)
 K_h\!\left(\widehat{\mathbf{p}}_u^{xy}(t)-\mathbf{c}_g\right),
 \label{eq:underserved-intensity}
\end{align}
where $\widehat{\mathbf{p}}_u^{xy}$ denotes the horizontal component
of the fused sensing estimate. The local underserved share is
\begin{equation}
 Z_g(t)=\frac{H_g(t)}{\rho_g(t)+\epsilon},
 \label{eq:underserved-share}
\end{equation}
where $\epsilon>0$ prevents division by zero. Using both $H_g$ and
$Z_g$ distinguishes a dense cluster containing many weak UEs from a
sparsely populated cell with a high underserved fraction.

\subsection{RF-Based Future Hotspot Prediction}

For proactive control, a Random Forest (RF) classifier operates on
one feature vector per grid cell and reporting epoch:
\begin{equation}
 \mathbf{x}_g(t)=
 [\rho_g,H_g,Z_g,\overline{R}_g,R_g^{\min},
 \overline{P}_g,\overline{C}_g,\overline{D}_g,
 \overline{L}_g,\overline{v}_g,\overline{\sigma}_g]^{T},
 \label{eq:rf-features}
\end{equation}
where the barred quantities are grid-level aggregates,
$\overline{v}_g$ is sensed UE speed, and
$\overline{\sigma}_g$ is mean localisation uncertainty. Features
whose measurements are unavailable are omitted.

For prediction horizon $\Delta_p$, let $N_g^\star$ and $Z_g^\star$
denote the oracle UE count and underserved share constructed from
simulation ground truth for training labels only. The future grid label is
\begin{equation}
 y_g(t)=\mathbf{1}\!\left\{
 N_g^\star(t+\Delta_p)\geq N_{\min}\ \land\
 Z_g^\star(t+\Delta_p)\geq\beta
 \right\},
 \label{eq:future-hotspot-label}
\end{equation}
where $N_{\min}$ prevents isolated UEs from forming a hotspot and
$\beta$ is the minimum underserved fraction. Ground-truth coordinates
are not included in $\mathbf{x}_g(t)$ and are unavailable during
inference. The RF estimates
\begin{equation}
 q_g(t)=
 \Pr\!\left[y_g(t)=1\mid\mathbf{x}_g(t)\right].
 \label{eq:rf-output}
\end{equation}
The predicted underserved region is therefore
\begin{equation}
 \widehat{\mathcal{G}}_{\mathrm{US}}(t)=
 \{g\in\mathcal{G}:q_g(t)\geq\tau\},
 \label{eq:predicted-region}
\end{equation}
where $\tau$ is selected using validation data. Independent
simulation seeds or trajectories are assigned to training,
validation, and testing to prevent leakage between temporally
adjacent samples.

\subsection{Hotspot Clustering and UAV Repositioning}

Let $K\leq|\mathcal{Q}|$ be the number of UAVs available for
repositioning. The target hotspot centroids are obtained through
probability-weighted K-means:
\begin{equation}
 \min_{\boldsymbol{\mu}_1,\ldots,\boldsymbol{\mu}_K}
 \sum_{g\in\widehat{\mathcal{G}}_{\mathrm{US}}(t)}
 q_g(t)H_g(t)
 \min_{k\in\{1,\ldots,K\}}
 \left\|\mathbf{c}_g-\boldsymbol{\mu}_k\right\|_2^2.
 \label{eq:weighted-kmeans}
\end{equation}
If $\mathbf{q}_j(t)$ is the horizontal position of UAV $j$, the
UAV-to-hotspot assignment minimises total repositioning distance:
\begin{align}
 \min_{a_{j,k}}\quad&
 \sum_{j\in\mathcal{Q}}\sum_{k=1}^{K}
 a_{j,k}\left\|\mathbf{q}_j(t)-\boldsymbol{\mu}_k(t)\right\|_2
 \label{eq:uav-assignment-objective}\\
 \mathrm{s.t.}\quad&
 \sum_k a_{j,k}\leq1,\quad
 \sum_j a_{j,k}=1,\quad
 a_{j,k}\in\{0,1\}.
 \label{eq:uav-assignment-constraints}
\end{align}
For the centroid $k(j)$ assigned to UAV $j$, its next position is
\begin{equation}
 \mathbf{q}_j(t+\Delta t)=\mathbf{q}_j(t)+
 \min\!\left\{v_{\max}\Delta t,
 \left\|\boldsymbol{\mu}_{k(j)}-\mathbf{q}_j(t)\right\|_2\right\}
 \frac{\boldsymbol{\mu}_{k(j)}-\mathbf{q}_j(t)}
 {\left\|\boldsymbol{\mu}_{k(j)}-\mathbf{q}_j(t)\right\|_2},
 \label{eq:uav-repositioning}
\end{equation}
subject to the configured flight-area and altitude constraints.
If the UAV is already at its assigned centroid, it retains its current
position and the final direction term in \eqref{eq:uav-repositioning}
is defined as zero.

\subsection{Evaluation Metrics}

Let $\mathcal{G}_{\mathrm{US}}(t)$ be an oracle underserved grid
constructed from true UE positions for evaluation only. Spatial
precision, recall, and F1 compare
$\widehat{\mathcal{G}}_{\mathrm{US}}(t)$ against
$\mathcal{G}_{\mathrm{US}}(t)$. After optimal one-to-one matching
$\pi(\cdot)$ between oracle centroids $\boldsymbol{\mu}_k$ and
estimated centroids $\widehat{\boldsymbol{\mu}}_k$, the centroid error
is
\begin{equation}
 E_{\mathrm{centroid}}=
 \frac{1}{K}\sum_{k=1}^{K}
 \left\|\boldsymbol{\mu}_k-
 \widehat{\boldsymbol{\mu}}_{\pi(k)}\right\|_2.
 \label{eq:centroid-error-final}
\end{equation}
The end-to-end control benefit is measured after UAV repositioning
using the underserved-UE fraction, coverage probability,
fifth-percentile throughput, packet-delivery ratio, and radio-outage
time. This separates sensing accuracy from the resulting improvement
in communication service.
